\clearpage\fancyhead[L]{\small JUNO.}\hypertarget{native-380-lower}{}\label{native-380-lower}
\begin{center}\textbf{JUNO.}\par\vspace{4pt}\hrule\vspace{4pt}\fontsize{10}{12}\selectfont Göttingische gelehrte Anzeigen. No. 19. p. 189. 31 January 1814.\par\vspace{4pt}\hrule\vspace{4pt}\end{center}

\par\indent As an addendum to our recent report on Juno (see no. 204 from the previous year), we include three observations of this planet by Professor \textsc{Bessel}:

\begin{center}\fontsize{10}{12}\selectfont\setlength{\tabcolsep}{3pt}\renewcommand{\arraystretch}{1.15}\begin{tabular}{l|c|rrr|c}
1812 & \shortstack{Mean time\\in Königsberg} & \multicolumn{3}{c|}{\shortstack{Right\\ascension}} & \shortstack{Southern\\declination}\\\hline
Nov. 18 & $12^{\mathrm u}\;12^{\mathrm m}\;26^{\mathrm s}5$ & $60^\circ$ & $41'$ & $21''5$ & ---\quad---\\
21 & $11\;58\;20.4$ & 60 & 6 & 42.5 & ---\quad---\\
Dec. 11 & $10\;25\;54.9$ & 56 & 39 & 0.8 & $3^\circ\;53'\;27''2$\\
\end{tabular}\end{center}

\par\indent These are the first fruits of the new Königsberg Observatory. Construction began three years ago, and Mr. \textsc{Bessel} had already moved into it last autumn. Astronomy owes the warmest thanks to the government for founding this institution so important to science, equipping it with excellent instruments, and completing it so quickly even in these years of hardship. Indeed, every month of delay in enabling an astronomer like \textsc{Bessel} to work to full effect would have been a loss to science.
